\section{A dozen combinatorial topoi}\label{sec:shapes}

The last batch is the generator at its most mechanical: take a shape used in combinatorics or computer science, check the Ore condition, read off the weak excluded middle. The checks are one line each; we found none of these examples in print except reflexive graphs \cite[Ex.~2.52]{hemelaer-rogers}.

\begin{theorem}\label{th:shapes}
\begin{center}\small
\begin{tabular}{llcc}
\toprule
structures & shape & excluded middle & weak excluded middle\\
\midrule
directed graphs & $V\rightrightarrows E$ & no & no\\
reflexive graphs & $V\rightrightarrows E$, with loops & no & no\\
simplicial sets & $\Delta$ & no & no\\
semi-simplicial sets & $\Delta_{\mathrm{inj}}$ & no & no\\
symmetric simplicial sets & non-empty finite sets & no & no\\
cubical sets & cube category & no & no\\
bisimplicial sets & $\Delta\times\Delta$ & no & no\\
augmented simplicial sets & $\Delta_+$ & no & \textbf{yes}\\
$\mathsf{FI}$-sets (finite sets, injections) & $\mathsf{FI}$ or $\mathsf{FI}^{\op}$ & no & \textbf{yes}\\
presheaves on finite sets & $\mathsf{Fin}$ & no & \textbf{yes}\\
the topos of trees & $\omega$ & no & \textbf{yes}\\
sets with an involution & $\B(\mathbb Z/2)$ & \textbf{yes} & \textbf{yes}\\
\bottomrule
\end{tabular}
\end{center}
\end{theorem}

\gen{} Bridges~\ref{br:bool} and~\ref{br:dm} on each shape. The \q{no} cases all fail Ore for the same reason: two distinct \q{vertex} arrows into the same object (the source and target of an edge, the two endpoints $[0]\rightrightarrows[1]$ of a segment) that no arrow can equalise, since every arrow into a vertex is constant. The \q{yes} cases have an initial object ($\Delta_+$, $\mathsf{FI}$, $\mathsf{Fin}$: complete any cospan by the empty object) or pushouts of injections ($\mathsf{FI}^{\op}$), or are a chain ($\omega$: complete by the minimum).

\direct{} For a \q{no}, the witness is always the same picture: in the one-edge graph (or $\Delta[1]$), let $A$ be one endpoint $u$; then $\neg A$ is the other endpoint $v$, $\neg\neg A = \{u\}$, and the edge is in neither, since it sees $A$ through one end and $\neg A$ through the other. For a \q{yes} with an initial object $\emptyset$: every element $x$ has a restriction $\res x{\emptyset}$ to the empty object, reached from every restriction of $x$; if any restriction of $x$ is in $A$, so is $\res x\emptyset$, and then every restriction of $x$ has a further restriction in $A$.\qed

\begin{example}[What the weak excluded middle says for augmented simplicial sets]
An augmented simplicial set is a simplicial set $X$ with a map $X_0\to X_{-1}$, a labelling of its components. For a sub-object $A$, a simplex lies in $\neg\neg A$ iff its label lies in $A_{-1}$, and in $\neg A$ iff it does not. The weak excluded middle is decided entirely by the augmentation.
\end{example}

\begin{remark}
For a programmer working with guarded recursion, the topos of trees is the standard model of the \q{later} modality \cite{birkedal}. That it is De Morgan means that its internal logic proves $\neg\varphi\vee\neg\neg\varphi$ for every proposition $\varphi$, although it does not prove $\varphi\vee\neg\varphi$; the proof above is two lines.
\end{remark}
